% ===================== فصل ۱ =====================
\chapter{شروعِ کار}
\label{ch:getting-started}

در این فصل نخستین گام‌ها را برمی‌داریم: کامپایلرِ سلام را آماده می‌کنیم،
نخستین برنامهٔ خود را می‌نویسیم و آن را اجرا می‌کنیم. نگران نباشید؛ اگر تا
امروز حتی یک خط برنامه ننوشته‌اید، این فصل دستِ شما را می‌گیرد و قدم‌به‌قدم
پیش می‌برد.

\section{برنامه‌نویسی یعنی چه؟}

رایانه ماشینی است بسیار سریع اما بسیار ساده‌دل. خودش نمی‌داند چه باید بکند؛
باید به او \emph{دستور} بدهیم. مجموعه‌ای از دستورها که به ترتیب نوشته شوند تا
رایانه کاری انجام دهد، یک \textbf{برنامه} نام دارد و زبانی که با آن این دستورها
را می‌نویسیم، یک \textbf{زبانِ برنامه‌نویسی} است.

سلام یکی از این زبان‌هاست؛ با این ویژگیِ ممتاز که می‌توانید دستورها را به زبانِ
\emph{فارسی} بنویسید. برای نمونه، جملهٔ «چیزی را چاپ کن» در سلام چنین نوشته
می‌شود:
\code{چاپ "سلام"}.

\begin{note}
رایانه دقیقاً همان کاری را می‌کند که به او می‌گویید نه بیشتر و نه کمتر. پس
اگر برنامه آن‌طور که می‌خواستید کار نکرد، احتمالاً اشتباه از دستوری است که داده‌اید،
نه از رایانه. این موضوع کاملاً طبیعی است و حتی برنامه‌نویسانِ حرفه‌ای هم هر روز
با آن روبه‌رو می‌شوند.
\end{note}

\section{کامپایلر چیست و چرا به آن نیاز داریم؟}

برنامه‌ای که شما می‌نویسید برای \emph{انسان} خواناست، اما رایانه آن را
مستقیماً نمی‌فهمد. به برنامه‌ای نیاز داریم که نوشتهٔ ما را به زبانی که رایانه
می‌فهمد ترجمه کند. این مترجم، \textbf{کامپایلر} نام دارد. کامپایلرِ زبانِ سلام
\code{salam} نامیده می‌شود.

\code{salam} کارهای زیادی برای ما انجام می‌دهد: متنِ برنامه را می‌خواند،
درستیِ آن را وارسی می‌کند، خطاها را گزارش می‌دهد و در پایان یک فایلِ اجرایی
می‌سازد که می‌توانید آن را روی رایانه اجرا کنید.

\begin{deep}[نگاهی به درون: مسیرِ یک برنامه]
وقتی \code{salam} برنامهٔ شما را می‌سازد، آن را از چند مرحله می‌گذراند:
\emph{تحلیلِ واژگانی} (شکستنِ متن به نشانه‌ها)، \emph{تجزیه} (ساختنِ درختِ
نحوی)، \emph{تحلیلِ معنایی} (وارسیِ انواع و نام‌ها)، و سرانجام \emph{تولیدِ کد}
که برنامهٔ شما را به زبانِ \lr{C} ترجمه می‌کند و سپس به فایلِ اجرایی تبدیل
می‌نماید. در سراسرِ این کتاب گاه‌به‌گاه به این مراحل اشاره خواهیم کرد، اما برای
نوشتنِ برنامه نیازی به دانستنِ جزئیاتِ آن‌ها نیست.
\end{deep}

\section{آماده‌سازیِ کامپایلر}

برای کار با سلام به کامپایلرِ \code{salam} نیاز دارید. این کامپایلر از روی
کدِ منبع ساخته می‌شود. در پوشهٔ اصلیِ پروژهٔ سلام، فرمانِ زیر را اجرا کنید:

\begin{salamen}[language={}]
sh tools/build-compiler.sh
\end{salamen}

پس از پایان، فایلِ اجرایی \code{salam} (یا در ویندوز \code{salam.exe}) در
پوشهٔ پروژه ساخته می‌شود. راهِ دیگر، ساخت با \lr{CMake} است:

\begin{salamen}[language={}]
cmake -B build
cmake --build build
\end{salamen}

\begin{tip}
برای آنکه مطمئن شوید کامپایلر درست ساخته شده است، فرمانِ زیر را اجرا کنید تا
شمارهٔ نسخه را ببینید:
\begin{salamen}[language={}]
salam -v
\end{salamen}
\end{tip}

\begin{note}
در این کتاب فرض می‌کنیم \code{salam} در دسترسِ شماست؛ یعنی یا در پوشهٔ جاری
قرار دارد (آن را با \code{./salam} صدا می‌زنید) یا در مسیرِ سیستم نصب شده است
(آن را با \code{salam} صدا می‌زنید). برای سادگی، در سراسرِ کتاب \code{salam}
می‌نویسیم.
\end{note}

\section{نخستین برنامه: سلام، دنیا!}

سنّتی دیرینه در میانِ برنامه‌نویسان هست: نخستین برنامه‌ای که در هر زبانِ تازه
می‌نویسند، برنامه‌ای است که عبارتِ «سلام، دنیا!» را نمایش می‌دهد. بیایید همین
سنّت را پاس بداریم.

فایلی به نامِ \code{01-hello.salam} بسازید و این چند خط را در آن بنویسید:

\begin{salam}
// تابع اصلی در واقع نقطه شروع نرم افزار شماست.
تابع اصلی:
    نام رشته = "سلام"
    چاپ "سلام،", نام
    چاپ "۲ + ۳ =", 2 + 3
پایان
\end{salam}

\begin{note}
پسوندِ همهٔ فایل‌های برنامهٔ سلام \code{.salam} است. این پسوند به ابزارها
می‌گوید که محتوای فایل، کدِ زبانِ سلام است.
\end{note}

\subsection{اجرای برنامه}

ساده‌ترین راه برای دیدنِ نتیجه، استفاده از مفسرِ سلام است. در ترمینال بنویسید:

\begin{salamen}[language={}]
salam exec 01-hello.salam
\end{salamen}

و خروجی چنین خواهد بود:

\begin{salamen}[language={}]
سلام، سلام
۲ + ۳ = 5
\end{salamen}

اگر می‌خواهید یک فایلِ اجراییِ مستقل بسازید (که سریع‌تر اجرا می‌شود و برای
انتشارِ برنامه مناسب است)، از فرمانِ \code{build} استفاده کنید:

\begin{salamen}[language={}]
salam build 01-hello.salam --output=hello.exe
./hello.exe
\end{salamen}

\begin{tip}
تفاوتِ \code{exec} و \code{build} در چیست؟ \code{exec} برنامه را
\emph{بی‌درنگ} اجرا می‌کند و برای آزمایش و یادگیری عالی است. \code{build} یک
فایلِ اجرایی می‌سازد که می‌توانید بارها و بدونِ نیاز به کامپایلر اجرایش کنید و
برای تحویلِ نهاییِ برنامه به‌کار می‌رود. در سراسرِ این کتاب، برای سادگی بیشتر
از \code{exec} استفاده می‌کنیم.
\end{tip}

\section{کالبدشکافیِ نخستین برنامه}

اکنون خط‌به‌خط ببینیم در این برنامه چه گذشت. هیچ بخشی را سرسری نگیرید؛ همین
چند خط، دروازهٔ ورود به همهٔ مفاهیمِ بعدی است.

\paragraph{خطِ نخست یک توضیح:}

هر خطی که با \code{//} آغاز شود، یک \textbf{توضیح} (comment) است و کامپایلر آن
را نادیده می‌گیرد. توضیحات برای \emph{انسان} نوشته می‌شوند، نه برای رایانه.
شکلِ ویژهٔ \code{//!} یک \emph{اطلاعاتِ سرآمد} (metadata) است که به ویرایشگرها
و محیط‌های آزمایش می‌گوید این فایل یک برنامهٔ اجرایی (\code{app}) با عنوانِ
«سلام» است. برای اجرای برنامه ضروری نیست، اما عادتی خوب است.

\paragraph{خطِ دوم آغازِ برنامه:}
\begin{salam}
تابع اصلی:
\end{salam}
این خط می‌گوید: «یک \textbf{تابعِ} به نامِ \textbf{اصلی} را تعریف کن.» تابع،
بسته‌ای از دستورهاست که نامی دارد. تابعی که نامش \code{اصلی} است، نقطهٔ
\emph{شروعِ} هر برنامهٔ سلام است؛ یعنی وقتی برنامه را اجرا می‌کنید، رایانه از
اینجا کار را آغاز می‌کند. نشانهٔ دونقطه (\code{:}) در پایانِ این خط می‌گوید که
بدنهٔ تابع از خطِ بعد آغاز می‌شود.

\begin{note}
نامِ \code{اصلی} در فارسی، دقیقاً معادلِ \code{main} در انگلیسی است. اگر همین
برنامه را به انگلیسی بنویسید، خطِ نخست چنین می‌شود: \code{func main:}.
\end{note}

\paragraph{خطِ سوم یک متغیر:}
\begin{salam}
    نام رشته = "سلام"
\end{salam}
اینجا یک \textbf{متغیر} به نامِ \code{نام} می‌سازیم. یک متغیر را مانندِ
جعبه‌ای تصور کنید که چیزی در آن نگه می‌داریم. در این جعبه، متنِ
\code{"سلام"} را گذاشته‌ایم. واژهٔ \code{رشته} می‌گوید که نوعِ این جعبه از
جنسِ \emph{متن} است (در فصلِ بعد بیشتر دربارهٔ انواع خواهیم گفت). هر متنی را
میانِ دو علامتِ نقل‌قول (\code{"}\,) می‌نویسیم.

دقت کنید که این خط با چند فاصله از سرِ سطر آغاز شده است. به این فاصله
\textbf{تورفتگی} (indentation) می‌گوییم و نشان می‌دهد که این خط \emph{درونِ}
تابعِ \code{اصلی} قرار دارد.

\paragraph{خطِ چهارم چاپِ خروجی:}
\begin{salam}
    چاپ "سلام،", نام
\end{salam}
واژهٔ \code{چاپ} به رایانه می‌گوید چیزی را روی صفحه نمایش بدهد. اینجا دو چیز
را چاپ می‌کنیم: نخست متنِ ثابتِ \code{"سلام،"} و سپس محتوای متغیرِ \code{نام}
(که همان \code{"سلام"} است). این دو با کاما (\code{,}) از هم جدا شده‌اند و سلام
میانِ آن‌ها یک فاصله می‌گذارد. به همین دلیل خروجی «سلام، سلام» می‌شود.

\paragraph{خطِ پنجم یک محاسبهٔ کوچک:}
\begin{salam}
    چاپ "۲ + ۳ =", 2 + 3
\end{salam}
سلام تنها متن چاپ نمی‌کند؛ می‌تواند حساب هم بکند. عبارتِ \code{2 + 3} پیش از
چاپ محاسبه می‌شود و نتیجه (یعنی \code{5}) چاپ می‌گردد. توجه کنید که
\code{"۲ + ۳ ="} درونِ نقل‌قول است و پس یک \emph{متنِ} ثابت است که بی‌کم‌وکاست
چاپ می‌شود، اما \code{2 + 3} بیرونِ نقل‌قول است و پس یک \emph{محاسبه} است.

\begin{warn}
به تفاوتِ ارقامِ فارسی (۲، ۳) و ارقامِ لاتین (\code{2}، \code{3}) دقت کنید.
درونِ یک متنِ نقل‌قول‌دار، هر رقمی فقط یک \emph{نویسه} است و چاپ می‌شود؛ اما
برای \emph{محاسبه}، سلام از ارقامِ لاتین استفاده می‌کند. پس \code{2 + 3} یک
جمعِ واقعی است، اما اگر بنویسید \code{"۲ + ۳"} تنها یک متن است و جمعی در کار
نیست.
\end{warn}

\paragraph{خطِ پایانی پایانِ تابع:}
\begin{salam}
پایان
\end{salam}
واژهٔ \code{پایان} به سلام می‌گوید که بدنهٔ تابعِ \code{اصلی} به پایان رسید.
هر تابعی که با \code{:} باز می‌شود، باید با \code{پایان} بسته شود. این قاعده‌ای
است که در سراسرِ زبان تکرار می‌شود: بلوک‌ها با دونقطه باز و با \code{پایان}
بسته می‌شوند.

\section{همان برنامه به انگلیسی}

یکی از زیبایی‌های سلام این است که هر برنامه را می‌توان به انگلیسی هم نوشت.
این برنامه دقیقاً معادلِ انگلیسیِ نمونهٔ بالاست:

\begin{salamen}
// Hello
func main:
    name str = "Salam"
    println "Hello,", name
    println "2 + 3 =", 2 + 3
end
\end{salamen}

می‌بینید که ساختار یکسان است؛ تنها واژگان تغییر کرده‌اند:
\code{تابع}~$\leftrightarrow$~\code{func}،
\code{اصلی}~$\leftrightarrow$~\code{main}،
\code{رشته}~$\leftrightarrow$~\code{str}،
\code{چاپ}~$\leftrightarrow$~\code{println} و
\code{پایان}~$\leftrightarrow$~\code{end}. جدولِ کاملِ این معادل‌ها را در
پیوستِ پایانِ کتاب خواهید یافت.

\begin{note}
خطِ \code{// LANG: fa} که در بسیاری از نمونه‌های این کتاب می‌بینید، به کامپایلر
می‌گوید زبانِ این فایل فارسی است. اگر این خط را ننویسید، سلام پیش‌فرض را زبانِ
انگلیسی می‌گیرد.
\end{note}

\section{چاپ با خط و بدونِ خط}

سلام دو شیوهٔ اصلی برای نمایشِ خروجی دارد:

\begin{itemize}
  \item \code{چاپ} (معادلِ \code{println}): چیزی را نمایش می‌دهد و سپس به
        \emph{خطِ بعد} می‌رود.
  \item \code{بنویس} (معادلِ \code{print}): چیزی را نمایش می‌دهد اما به خطِ
        بعد \emph{نمی‌رود}؛ یعنی نوشتهٔ بعدی در همان خط ادامه می‌یابد.
\end{itemize}

به این مثال دقت کنید:

\begin{salam}
تابع اصلی:
    بنویس "سلام "
    بنویس "دنیا"
    چاپ ""
    چاپ "خطِ تازه"
پایان
\end{salam}

خروجی:

\begin{salamen}[language={}]
سلام دنیا
خطِ تازه
\end{salamen}

دو فرمانِ \code{بنویس} هر دو در یک خط نوشتند و سپس \code{چاپ ""} (که یک متنِ
خالی را چاپ می‌کند) ما را به خطِ بعد بُرد. این ترفندِ ساده نوشتن با
\code{بنویس} و رفتن به خطِ بعد با \code{چاپ ""} را در فصلِ حلقه‌ها برای
کشیدنِ شکل‌ها بارها به‌کار خواهیم برد.

\section{خلاصهٔ فصل}

در این فصل آموختیم:
\begin{itemize}
  \item برنامه مجموعه‌ای از دستورهاست و کامپایلرِ \code{salam} آن را برای
        رایانه ترجمه می‌کند.
  \item هر برنامهٔ سلام از تابعِ \code{اصلی} آغاز می‌شود.
  \item بلوک‌ها با \code{:} باز و با \code{پایان} بسته می‌شوند.
  \item با \code{چاپ} و \code{بنویس} خروجی نمایش می‌دهیم.
  \item توضیحات با \code{//} نوشته می‌شوند و کامپایلر آن‌ها را نادیده می‌گیرد.
  \item هر برنامه را می‌توان به فارسی یا انگلیسی نوشت.
\end{itemize}

\section{تمرین‌ها}

\exercise برنامه‌ای بنویسید که نام و سنِ شما را در دو خطِ جداگانه چاپ کند.

\exercise برنامه‌ای بنویسید که حاصلِ \code{12 * 8} را به همراهِ یک متنِ توضیحی
چاپ کند؛ برای نمونه: \lr{\texttt{12 * 8 = 96}}.

\exercise برنامهٔ «سلام، دنیا!» را یک بار به فارسی و یک بار به انگلیسی بنویسید
و هر دو را اجرا کنید. آیا خروجی یکسان است؟

\exercise با استفاده از \code{بنویس}، نام و نام‌خانوادگیِ خود را در \emph{یک}
خط و با یک فاصله میانشان چاپ کنید.
